\documentclass[11pt,a4paper]{article}
\usepackage[utf8]{inputenc}
\usepackage[spanish]{babel}
\usepackage{geometry}
\geometry{a4paper, margin=1in}
\usepackage{hyperref}
\usepackage{booktabs}
\usepackage{graphicx}
\usepackage{xcolor}
\usepackage{titlesec}
\usepackage{array}

\hypersetup{
    colorlinks=true,
    linkcolor=blue,
    filecolor=magenta,      
    urlcolor=cyan,
    pdftitle={Manual de Uso: Anuncios y Pop-ups - ECG Corporativo},
    pdfpagemode=FullScreen,
}

\titleformat{\section}
  {\normalfont\Large\bfseries\color{blue!80!black}}
  {\thesection}{1em}{}

\titleformat{\subsection}
  {\normalfont\large\bfseries\color{blue!60!black}}
  {\thesubsection}{1em}{}

\title{\textbf{Manual de Uso: Gestión de Anuncios y Pop-ups\\ ECG Corporativo}}
\author{\textbf{Equipo de Desarrollo y Soporte}}
\date{\today}

\begin{document}

\maketitle
\tableofcontents
\newpage

\section{Introducción al Sistema de Anuncios}
El sistema de anuncios permite a los administradores y trabajadores mostrar ventanas emergentes (\textit{pop-ups}) al inicio de la sesión del usuario. Estos anuncios pueden configurarse de manera global (para el menú principal del portal) o segmentarse específicamente para cada una de las empresas de la corporación.

\subsection{Características Clave}
\begin{itemize}
    \item \textbf{Segmentación por destino:} Permite elegir si el anuncio aparece en el Portal Principal o en una empresa específica de la corporación.
    \item \textbf{Control de vigencia:} Configuración de una fecha límite; el anuncio dejará de mostrarse automáticamente después de esa fecha.
    \item \textbf{Contador regresivo (cuenta atrás):} Los pop-ups calculan y muestran el tiempo restante exacto de vigencia de forma dinámica y visual.
    \item \textbf{Habilitación rápida:} Permite activar o desactiva anuncios con un solo interruptor sin necesidad de eliminarlos de la base de datos.
\end{itemize}

\section{Formatos Disponibles}
El sistema soporta dos estilos de diseño altamente optimizados y adaptables:

\subsection{Modo Estándar (Tarjeta con Texto)}
Ideal para avisos informativos, novedades o promociones detalladas.
\begin{itemize}
    \item \textbf{Elementos en pantalla:} Ícono representativo, etiqueta de tipo (\textit{badge}), título, subtítulo, descripción textual completa, imagen descriptiva (opcional) y un botón de llamada a la acción (\textit{CTA}) personalizable.
    \item \textbf{Uso recomendado:} Comunicados corporativos, avisos de mantenimiento, lanzamiento de nuevos servicios o promociones que requieran explicación.
\end{itemize}

\subsection{Modo Solo Imagen (Banner con Cuenta Atrás)}
Diseñado específicamente para campañas de marketing visual u ofertas de alto impacto gráfico.
\begin{itemize}
    \item \textbf{Elementos en pantalla:} Únicamente la imagen de tu banner, una etiqueta opcional en la esquina superior izquierda, y una \textbf{barra inferior dedicada} para la cuenta atrás.
    \item \textbf{Visualización optimizada:} El pop-up se amplía automáticamente a un formato más grande (\textit{max-w-xl}) y el temporizador se coloca debajo de la foto. De esta manera, \textbf{los números nunca tapan los textos o caras} de la imagen.
    \item \textbf{Interactividad:} Si se configura un enlace, \textbf{toda la imagen se vuelve cliqueable}, actuando como un acceso directo a la promoción (por ejemplo, a un chat de WhatsApp o landing page).
    \item \textbf{Uso recomendado:} Flyers de descuentos, ofertas especiales de fin de semana, o posters publicitarios prediseñados de gran calidad.
\end{itemize}

\newpage

\section{Guía de Configuración de Campos (Formulario)}
Al crear o editar un anuncio en el panel administrativo, se deben configurar los siguientes parámetros:

\begin{table}[htbp]
\centering
\caption{Campos de Configuración de Anuncios}
\label{tab:campos}
\begin{tabular}{m{3cm} c m{9cm}}
\toprule
\textbf{Campo} & \textbf{Obligatorio} & \textbf{Descripción / Recomendación} \\
\midrule
\textbf{Tipo} & Sí & Define el color del pop-up y la etiqueta del tipo. Opciones: \textit{Oferta}, \textit{Novedad}, \textit{Evento}, \textit{Aviso} o \textit{Promoción}. \\
\addlinespace
\textbf{Destino} & Sí & Lugar de despliegue del anuncio. Opciones: \textit{Portal Principal} o la empresa correspondiente. \\
\addlinespace
\textbf{Modo Solo Imagen} & Opcional & Interruptor que cambia el diseño al formato visual de solo imagen (oculta los campos de texto estándar). \\
\addlinespace
\textbf{URL de la Imagen} & Sí (en Solo Imagen) & Enlace directo a la imagen. Debe ser una URL de internet pública (ej. \url{https://servidor.com/imagen.png}). \\
\addlinespace
\textbf{Título} & Sí (en Estándar) & Título llamativo del anuncio (se autodefine en el servidor si usas Solo Imagen). \\
\addlinespace
\textbf{Descripción} & Sí (en Estándar) & Cuerpo del mensaje o términos de la promoción (se autodefine en el servidor si usas Solo Imagen). \\
\addlinespace
\textbf{Badge} & Opcional & Pequeño texto destacado en fondo amarillo. Ej: \texttt{¡GRATIS!}, \texttt{30\% OFF}, \texttt{EXCLUSIVO}. \\
\addlinespace
\textbf{Link al hacer clic} & Opcional & URL a donde se redirigirá al usuario cuando haga clic en la imagen (Solo Imagen) o en el botón (Estándar). Ej. \url{https://wa.me/521XXXXXXXXXX} \\
\addlinespace
\textbf{Fecha de vencimiento} & Sí & Fecha límite de la campaña. Se utiliza para calcular el reloj de cuenta atrás. El pop-up expira y desaparece a las 23:59:59 de ese día. \\
\addlinespace
\textbf{Publicar ya} & Sí & Interruptor de estado del anuncio. Si se apaga, el anuncio se guarda como borrador inactivo. \\
\bottomrule
\end{tabular}
\end{table}

\section{Gestión en el Panel Administrativo}

\subsection{Crear un anuncio}
\begin{enumerate}
    \item Dirígete a la sección de \textbf{Gestión de Anuncios} en el panel de administración.
    \item Haz clic en el botón \textbf{`+ Nuevo anuncio'} en la esquina superior derecha.
    \item Rellena los campos correspondientes según el formato deseado y presiona \textbf{`Crear anuncio'}.
\end{enumerate}

\subsection{Modificar un anuncio}
\begin{itemize}
    \item \textbf{Desactivar temporalmente:} Haz clic en el interruptor de estado (icono de interruptor verde/gris) en la lista para ocultar el anuncio instantáneamente sin borrarlo.
    \item \textbf{Editar detalles:} Haz clic en el botón azul de editar (icono de lápiz 📝). Modifica los campos necesarios y guarda los cambios.
    \item \textbf{Eliminar permanentemente:} Haz clic en el botón rojo de eliminar (icono de papelera 🗑️) y confirma la acción (disponible únicamente para nivel Administrador o superior).
    \item \textbf{Vista Previa en tiempo real:} Utiliza el botón de ojo (👁️) en la lista para observar exactamente cómo se mostrará el pop-up a tus usuarios finales sin necesidad de guardarlo primero.
\end{itemize}

\newpage

\section{Referencia Técnica (Base de Datos Supabase)}
El sistema almacena la configuración de cada pop-up en la tabla \texttt{anuncios} en PostgreSQL. La estructura SQL de la base de datos es la siguiente:

\begin{verbatim}
CREATE TABLE IF NOT EXISTS anuncios (
  id            UUID          PRIMARY KEY DEFAULT gen_random_uuid(),
  titulo        TEXT          NOT NULL,
  subtitulo     TEXT          DEFAULT '',
  cuerpo        TEXT          NOT NULL,
  tipo          TEXT          NOT NULL DEFAULT 'aviso'
                                CHECK (tipo IN ('oferta','novedad','evento',
                                                'aviso','promocion')),
  icono         TEXT          DEFAULT 'Bell'
                                CHECK (icono IN ('Tag','Zap','Gift',
                                                'Bell','Sparkles')),
  badge         TEXT          DEFAULT '',
  destino       TEXT          NOT NULL DEFAULT 'portal'
                                CHECK (destino IN ('portal','empresa_1',
                                                'empresa_2','empresa_3')),
  cta_texto     TEXT          DEFAULT '',
  cta_link      TEXT          DEFAULT '',
  imagen_url    TEXT          DEFAULT '',
  solo_imagen   BOOLEAN       NOT NULL DEFAULT FALSE,
  fecha_fin     DATE          DEFAULT NULL,
  activo        BOOLEAN       NOT NULL DEFAULT TRUE,
  creado_por    TEXT          NOT NULL DEFAULT '',
  usuario_id    INTEGER,
  created_at    TIMESTAMPTZ   NOT NULL DEFAULT now(),
  updated_at    TIMESTAMPTZ   NOT NULL DEFAULT now()
);
\end{verbatim}

\end{document}
